\section{空间解析几何 II}

\subsection{空间平面与直线的夹角}

\subsubsection{直线与直线的夹角}

直线之间的夹角即为方向向量之间的夹角。设两直线 $\ell_1, \ell_2$ 的方向向量为 $\vect{u}_1, \vect{u}_2$，那么：
$$
\cos \theta = \cos \langle \vect{u}_1, \vect{u}_2 \rangle = \frac{\vect{u}_1 \vdot \vect{u}_2}{\norm{\vect{u}_1} \cdot \norm{\vect{u}_2}}
$$
因此：
$$
\theta = \arccos \frac{\abs{\vect{u}_1 \vdot \vect{u}_2}}{\norm{\vect{u}_1} \cdot \norm{\vect{u}_2}}
$$

\subsubsection{直线与平面的夹角}

直线与平面的夹角为方向向量与法向量之间夹角的余角。设直线 $\ell$ 方向向量为 $\vect{u}$，平面 $\alpha$ 法向量为 $\vect{n}$，那么：
$$
\sin \theta = \cos \langle \vect{u}, \vect{n} \rangle = \frac{\vect{u} \vdot \vect{n}}{\norm{\vect{u}} \cdot \norm{\vect{n}}}
$$
因此
$$
\theta = \arcsin \frac{\abs{\vect{u} \vdot \vect{n}}}{\norm{\vect{u}} \cdot \norm{\vect{n}}}
$$

\subsubsection{平面与平面的夹角}

平面之间的夹角就是法向量之间的夹角。设平面 $\alpha_1, \alpha_2$ 的法向量为 $\vect{n}_1, \vect{n}_2$，那么：
$$
\cos \theta = \cos \langle \vect{n}_1, \vect{n}_2 \rangle = \frac{\vect{n}_1 \vdot \vect{n}_2}{\norm{\vect{n}_1} \cdot \norm{\vect{n}_2}}
$$
因此
$$
\cos = \arccos \frac{\abs{\vect{n}_1 \vdot \vect{n}_2}}{\norm{\vect{n}_1} \cdot \norm{\vect{n}_2}}
$$

\subsection{作业}

\begin{problem}
	习题 8.4.9
	\begin{solution}
		直线的方向向量为
		$$
		\vect{u} = (1, 1, 3) \cp (1, -1, -1) = (2, 4, -2)
		$$
		平面的法向量为
		$$
		\vect{n} = (1, -1, -1)
		$$
		那么可得
		$$
		\sin \theta = \frac{\vect{u} \vdot \vect{n}}{\norm{\vect{u}} \cdot \norm{\vect{n}}} = 0
		$$
		因此夹角 $\theta = 0$。
	\end{solution}
\end{problem}

\begin{problem}
	习题 8.4.11
	\begin{solution}
		两直线的方向向量为
		$$
		\begin{gathered}
			\vect{u}_1 = (1, 2, -1) \cp (1, -1, 1) = (1, -2, -3) \\
			\vect{u}_2 = (2, -1, 1) \cp (1, -1, 1) = (0, -1, -1)
		\end{gathered}
		$$
		因此，平面的法向量为
		$$
		\vect{n} = \vect{u}_1 \cp \vect{u}_2 = (-1, 1, -1)
		$$
		可得平面方程为
		$$
		-x + y - z = 0
		$$
	\end{solution}
\end{problem}

\begin{problem}
	习题 8.4.14
	\begin{solution}
		设 $M M_0$ 与 $\vect{s}$ 的夹角为 $\theta$。那么可知：
		$$
		d = \abs{M M_0} \sin \theta = \abs{M M_0} \cdot \frac{\norm{\overrightarrow{M M_0} \cp \vect{s}}}{\abs{M M_0} \cdot \norm{\vect{s}}} = \frac{\norm{\overrightarrow{M M_0} \cp \vect{s}}}{\norm{\vect{s}}}
		$$
	\end{solution}
\end{problem}

\begin{problem}
    习题 8.4.15
    \begin{solution}
        取直线上两点 $P_1(0, -1, -4), P_2(\frac{9}{7}, 0, -\frac{18}{7})$，其分别投影到平面 $4x - y + z = 1$ 上的结果为：
		$$
		P_1'\ab(\frac{8}{9}, -\frac{11}{9}, -\frac{34}{9}), P_2'\ab(\frac{59}{63}, \frac{11}{126}, -\frac{335}{126})
		$$
		于是直线方程为：
		$$
		(x, y, z) = \ab(\frac{8}{9}, -\frac{11}{9}, -\frac{34}{9}) + \lambda (2, 110, 94)
		$$
    \end{solution}
\end{problem}